\section{New mod 10 identities}
\label{sec:mod10}

Varying the linear terms in the exponent of $q$ from~\eqref{eqn:mod10_intro}, we have the following list of identities where the products are related to the
principal characters standard (i.e., highest-weight, integrable) $\mathrm{D}_4^{(3)}$ modules of level $4$.
For notation regarding $\mathrm{D}_4^{(3)}$, see Carter's book \cite[P.\ 608]{Car-book}. 
For more information on principal characters and computational techniques for them, see the work of Bos~\cite{Bos}.
\begin{align}
&\sum_{i,j,k,\ell \ge 0}
\frac{q^{2i^2 + 6 ij + 4ik + 2i\ell + 6j^2 + 8 jk + 4 j\ell + 3k^2+3k\ell+\ell^2}}
{\left(q^2;q^2\right)_i\left(q^2;q^2\right)_j\left(q;q\right)_k\left(q;q\right)_\ell}\nonumber\\
&\quad=
\dfrac{1}{\theta(q,q^2,q^4;q^{10})} 
=\chi(\Omega(2\Lambda_0+\Lambda_1)),
\label{eq:mod10A}
\\
&\sum_{i,j,k,\ell \ge 0}
\frac{q^{2i^2 + 6 ij + 4ik + 2i\ell + 6j^2 + 8 jk + 4 j\ell + 3k^2+3k\ell+\ell^2+2j+k+\ell}}
{\left(q^2;q^2\right)_i\left(q^2;q^2\right)_j\left(q;q\right)_k\left(q;q\right)_\ell}
\nonumber\\ 
&\quad=
\dfrac{1}{\theta(q^2,q^2,q^3;q^{10})}  
=\chi(\Omega(2\Lambda_1)),
\\
&\sum_{i,j,k,\ell \ge 0}
\frac{q^{2i^2 + 6 ij + 4ik + 2i\ell + 6j^2 + 8 jk + 4 j\ell + 3k^2+3k\ell+\ell^2+2i+2j+2k}}
{\left(q^2;q^2\right)_i\left(q^2;q^2\right)_j\left(q;q\right)_k\left(q;q\right)_\ell}
\nonumber\\ 
&\quad= 
\dfrac{1}{\theta(q,q^4,q^4;q^{10})}  
=\chi(\Omega(\Lambda_0+\Lambda_2)),
\\
&\sum_{i,j,k,\ell \ge 0}
\frac{q^{2i^2 + 6 ij + 4ik + 2i\ell + 6j^2 + 8 jk + 4 j\ell + 3k^2+3k\ell+\ell^2+2i+4j+3k+\ell}}
{\left(q^2;q^2\right)_i\left(q^2;q^2\right)_j\left(q;q\right)_k\left(q;q\right)_\ell}
\nonumber\\ 
&\quad=
\dfrac{1}{\theta(q^2,q^3,q^4;q^{10})}   
= \chi(\Omega(4\Lambda_0)). \label{eq:mod10D}
\end{align}
Recall that  identities for principal characters of  standard modules
of $\mathrm{D}_4^{(3)}$ with level $3$ were previously conjectured by the second and third authors~\cite{KanRus-idf}, and
analytic forms of these conjectures were found by Kur\c{s}ung\"{o}z~\cite{Kur-KR}. See also the thesis of the third author~\cite{Rus-thesis} for further
related identities.

We will devote the remainder of this section to a proof of these conjectures.

As in the previous section, we begin by introducing refinements of the quadruple sums, along with arbitrary linear terms in the exponent of $q$. However, this time, we use two variables, $x$ and $y$.
\begin{align}
\label{eqn:defR}
&R_{A,B,C,D}(x,y,q)\\
&= \sum_{i,j,k,\ell \ge 0}
\frac{x ^{2j + k + \ell } y ^{ i+j+k}q^{2i^2 + 6 ij + 4ik + 2i\ell + 6j^2 + 8 jk + 4 j\ell + 3k^2+3k\ell+\ell^2+Ai+Bj+Ck+D\ell}}
{\left(q^2;q^2\right)_i\left(q^2;q^2\right)_j\left(q;q\right)_k\left(q;q\right)_\ell}. \nonumber
\end{align}
As usual, we will drop the arguments $x,y,q$ when they are clear.

We have the following shifts of $R_{A,B,C,D}$:
\begin{align}
R_{A,B,C,D}(xq,y,q)&=R_{A,B+2,C+1,D+1}(x,y,q) \label{eqn:shiftRx},\\
R_{A,B,C,D}(x,yq,q)&=R_{A+1,B+1,C+1,D}(x,y,q) \label{eqn:shiftRy}.
\end{align}


The fundamental relations governing these sums are deduced easily.
\begin{prop}
\label{prop:mod10rels}
\begin{alignat}{3}
m_1(A,B,C,D):&\quad R_{A,B,C,D}-R_{A+2,B,C,D}-yq^{2+A}R_{A+4,B+6,C+4,D+2}&&=0 \label{eqn:relm1},\\
m_2(A,B,C,D):&\quad R_{A,B,C,D}-R_{A,B+2,C,D}-x^2yq^{6+B}R_{A+6,B+12,C+8,D+4}&&=0 \label{eqn:relm2},\\
m_3(A,B,C,D):&\quad R_{A,B,C,D}-R_{A,B,C+1,D}-xyq^{3+C}R_{A+4,B+8,C+6,D+3}&&=0 \label{eqn:relm3},\\
m_4(A,B,C,D):&\quad R_{A,B,C,D}-R_{A,B,C,D+1}-xq^{1+D}R_{A+2,B+4,C+3,D+2}&&=0 \label{eqn:relm4}.
\end{alignat}
\end{prop}


\begin{thm}
The series $R_{0,0,0,0}$ is the unique solution in $\ZZ[[x,y,q]]$ satisfying
the following:
\begin{align}
F(x,y,q) &= F(xq,y,q) + xqF\left(xq^2,y,q\right) \label{eqn:m10d1},\\
F(x,y,0) &= 1,\quad 
F(x,0,q) = \sum_{\ell\geq 0}\dfrac{q^{\ell^2}x^\ell}{(q)_\ell},\quad 
F(0,y,q) = \sum_{i\geq 0}\dfrac{q^{2i^2}y^i}{(q^2;q^2)_i}. \label{eqn:m10init}
\end{align}
\end{thm}
\begin{proof}
$R_{0,0,0,0}$ is seen to satisfy \eqref{eqn:m10init} easily.
Due to the shifts \eqref{eqn:shiftRx} and \eqref{eqn:shiftRy}, showing that $R_{0,0,0,0}$ satisfies
\eqref{eqn:m10d1} is equivalent to showing:
\begin{align}
R_{0,0,0,0} - R_{0,2,1,1} - xqR_{0,4,2,2} = 0.
\end{align}
This relation can be obtained as a linear combination of the fundamental relations \eqref{eqn:relm1}--\eqref{eqn:relm4}:
\begin{align}
-&m_1(-2,0,0,0)+m_1(-2,0,0,1)-xq\cdot m_1(0,4,2,2)
\nonumber\\
&+xq\cdot m_1(0,4,3,2)+m_2(0,0,0,1)+m_3(0,2,0,1)
\nonumber \\
&-xq\cdot m_3(2,4,2,2)+m_4(-2,0,0,0)-y\cdot m_4(2,6,4,2).
\end{align}

Now we prove the uniqueness. 
Suppose that $F(x,y,q)=\sum_{i,j,k\geq 0}f_{i,j,k}x^iy^jq^k$ is a solution to this system.
Note that for all $i,j,k\geq 0$, $f_{0,j,k}$, $f_{i,0,k}$ and $f_{i,j,0}$ 
are uniquely determined due to the initial conditions 
\eqref{eqn:m10init}.
For convenience, we define $f_{i,j,k}=0$ whenever any one or more of $i,j,k$ are negative.
\eqref{eqn:m10d1} translates to:
\begin{align}
f_{i,j,k}&=f_{i,j,k-i}+f_{i-1,j,k-2i+1}.
\label{eqn:frel}
\end{align}
Now we induct on $N=i+k$ and show that all $f_{i,j,k}$ are uniquely determined.
The case~$N=0$ follows easily since $f_{0,j,0}$ are uniquely known due to \eqref{eqn:m10init}.
Suppose that for all triples $(i,j,k)$ with $i+k<N$, the values $f_{i,j,k}$ are determined.
Pick a triple~$(I,J,K)$ such that $I+K=N$. 
If $I=0$, we know $f_{0,J,N}$ by \eqref{eqn:m10init}.
So, suppose that~$I\geq 1$. 
Then, the RHS of \eqref{eqn:frel} involves two terms:
\begin{enumerate}
\item $f_{I,J,K-I}$ for which $I+(K-I) = K < N$, (since $K+I=N$ and $I\geq 1$)
\item $f_{I-1,J,K-2I+1}$ for which $(I-1)+(K-2I+1) = K - I < N$ (since $K+I=N$ and $I\geq 1$).
\end{enumerate}
This means that the RHS of \eqref{eqn:frel} has been determined already, and this determines
the LHS uniquely.
\end{proof}

\begin{thm}
Identities \eqref{eq:mod10A}--\eqref{eq:mod10D} are true.
\end{thm}
\begin{proof}
It is easy to see using chapter 7 of Andrews' text~\cite{And-book} that 
the series  
\begin{align}
\left(\sum_{i\geq 0} \dfrac{x^iq^{i^2}}{(q)_i}\right)
\left(\sum_{j\geq 0} \dfrac{y^jq^{2j^2}}{(q^2;q^2)_j}\right)
\end{align}
satisfies \eqref{eqn:m10d1} and \eqref{eqn:m10init}.
Thus, due to uniqueness, we have:
\begin{align}
R_{0,0,0,0}(x,y,q)=\left(\sum_{i\geq 0} \dfrac{x^iq^{i^2}}{(q)_i}\right)
\left(\sum_{j\geq 0} \dfrac{y^jq^{2j^2}}{(q^2;q^2)_j}\right).
\end{align}
By the Rogers--Ramanujan identities \cite[Ch.\ 7]{And-book}, we now have:
\begin{align}
R_{0,0,0,0}(1,1,q)=\dfrac{1}{\theta(q;q^5)\theta(q^2;q^{10})},\\
R_{0,0,0,0}(q,1,q)=\dfrac{1}{\theta(q^2;q^5)\theta(q^2;q^{10})},\\
R_{0,0,0,0}(1,q^2,q)=\dfrac{1}{\theta(q;q^5)\theta(q^4;q^{10})},\\
R_{0,0,0,0}(q,q^2,q)=\dfrac{1}{\theta(q^2;q^5)\theta(q^4;q^{10})}.
\end{align}
This immediately proves~\eqref{eq:mod10A}--\eqref{eq:mod10D}.
\end{proof}
